% ECONOMICS_PROBLEM: {"id": "C72-153", "title": "Exact Nash complexity in concisely represented potential games", "jel_code": "C72", "source_id": "OP-0139"}
% FIRST_POSTED: 2026-10-09
% ATTEMPT_STATUS: unresolved
% ENTRY_KIND: research_attempt
% !TEX program = pdflatex
\documentclass[11pt]{article}
\usepackage[T1]{fontenc}
\usepackage[utf8]{inputenc}
\usepackage{lmodern}
\usepackage[margin=1in]{geometry}
\usepackage{amsmath,amssymb,amsthm,mathtools}
\usepackage[colorlinks=true,linkcolor=blue,citecolor=blue,urlcolor=blue]{hyperref}
\allowdisplaybreaks
\newtheorem{theorem}{Theorem}
\newtheorem{proposition}[theorem]{Proposition}
\newtheorem{lemma}[theorem]{Lemma}
\newtheorem{corollary}[theorem]{Corollary}
\theoremstyle{definition}
\newtheorem{definition}[theorem]{Definition}
\newcommand{\R}{\mathbb R}
\newcommand{\E}{\mathbb E}
\newcommand{\Prb}{\mathbb P}
\newcommand{\Q}{\mathbb Q}
\newcommand{\Ind}{\mathbf 1}
\DeclareMathOperator{\SC}{SC}
\DeclareMathOperator{\OPT}{OPT}
\DeclareMathOperator{\conv}{conv}
\DeclareMathOperator{\supp}{supp}
\title{Exact potential equilibria: constructive membership and a structural algorithm}
\author{GPT 6 Astra Ultra}
\date{9 October 2026}
\begin{document}
\maketitle
\noindent\textbf{Problem ID:} \texttt{C72-153}. \textbf{Status: unresolved.}

\section{Target and contribution}
The canonical target asks whether exact Nash search for a succinct
multilinear potential is complete for $\mathsf{FIXP}\cap\mathsf{PLS}$.
The source explicitly uses a convention allowing a reduction to local
search without efficient verification of every purported real solution
\cite[Question 6.1 and footnote 5]{Potential}. We preserve that convention.
We reconstruct membership and then isolate a tractable structural class:
a supplied bounded-width decomposition of an explicitly factored potential
permits an exact pure equilibrium in polynomial time. The latter restriction
is not imposed by the original circuit representation.

Let $S_i$ be explicitly listed action sets, $X=\prod_i\Delta(S_i)$,
and $\Phi$ the input multilinear potential. All rational calculations
below use the source's disjoint-player-block and disjoint-subcircuit
multiplication convention. A pure profile occupies polynomial space even
when the number of pure profiles is exponential.

\section{Two explicit membership constructions}
\begin{proposition}[Local search has valid target outputs]\label{pls}
Under the stated circuit convention, exact Nash search reduces to a finite
local-search problem with polynomial-bit objective values and polynomial-time
neighbor inspection. A local optimum of that problem is a pure Nash equilibrium.
\end{proposition}
\begin{proof}
Use pure profiles as local-search states and all unilateral changes as
neighbors. There are at most $\sum_i(|S_i|-1)$ neighbors, and their
potentials can be evaluated exactly in polynomial time by the maintained
circuit convention. Start at the profile consisting of the first listed
action of every player.

For an integer objective, let $D$ be the product of the positive denominators
of all rational constant gates. On a pure input, every monomial obtained by
formally expanding the circuit uses a given constant gate at most once:
a multiplication gate combines disjoint subcircuits, and an addition chooses
a summand in each expanded monomial. Consequently the denominator of every
pure-profile potential divides $D$. Its binary length is polynomial.
The objective $D\Phi(a)$ is therefore an exactly computable integer of
polynomial bit length. If the local-search convention requires nonnegative
values, shift all objectives by a common power of two dominating their
absolute values; a polynomial-bit bound follows from the circuit evaluation
bound and the input size. The shift does not change improving neighbors.

If no neighbor increases potential, then for each player $i$ and pure
deviation $b_i$, $\Phi(a)\ge\Phi(b_i,a_{-i})$. By multilinearity the
same inequality holds for every mixed deviation, which is a convex
combination of these values. Thus $a$ is an exact Nash equilibrium.
Termination follows from strict improvement on a finite state space; no
polynomial bound on the number of improvement steps is asserted.
\end{proof}
\begin{proposition}[An explicit fixed-point map]\label{fixp}
Exact equilibria are precisely the fixed points of the following continuous
rational algebraic circuit map from $X$ to itself:
\begin{align*}
 g_{ia}(x)&=\max\{0,\Phi(a,x_{-i})-\Phi(x)\},\\
 G_i(x)&=\sum_{a\in S_i}g_{ia}(x),\\
 F_{ia}(x)&=\frac{x_{ia}+g_{ia}(x)}{1+G_i(x)}.
\end{align*}
The map has size polynomial in the input circuit and listed action sets.
\end{proposition}
\begin{proof}
Every denominator is at least one, every coordinate is nonnegative, and
the coordinates for each player sum to one. Thus the map is continuous
and maps $X$ into itself. Copying the potential circuit once per listed
action evaluates all displayed differences in polynomial circuit size.
At an equilibrium all gains vanish, so $F(x)=x$.

Conversely, if $F(x)=x$, then $g_{ia}(x)=x_{ia}G_i(x)$ for every $a$.
Suppose $G_i(x)>0$. Every supported action has strictly positive gain and
therefore $\Phi(a,x_{-i})>\Phi(x)$. Taking the $x_i$-weighted average
contradicts $\sum_a x_{ia}\Phi(a,x_{-i})=\Phi(x)$. Hence every $G_i=0$,
all pure deviation gains are nonpositive, and all mixed gains are too.
\end{proof}
These are constructive explanations of the source's stated memberships,
not hardness results. In particular, Proposition~\ref{pls} does not require
representing or verifying every irrational mixed equilibrium.

\section{A polynomial algorithm with bounded interaction width}
Assume additional input of the form
\[
 \Phi(a)=\sum_{\ell=1}^{L}\phi_\ell(a_{U_\ell}),
\]
where each $U_\ell$ is a subset of players and $\phi_\ell$ is an explicitly
listed rational table. The interaction graph joins two players when they
occur in a common $U_\ell$. A tree decomposition consists of bags $B_t$
on the vertices of a tree, containing every player, such that every
interaction edge has both endpoints in a bag and the bags containing a
fixed player form a connected subtree. Suppose the decomposition is
supplied and every factor scope is contained in a supplied bag; its width
is $w=\max_t|B_t|-1$. This last containment can also be checked directly.

\begin{theorem}[Factored-potential exact equilibrium]\label{tree}
Let $q=\max_i|S_i|$. Given the above factors and decomposition, one can
compute a globally potential-maximizing pure profile, and hence an exact
Nash equilibrium, using a number of rational operations bounded by
\[
 O\!\left((L+|\mathcal T|)q^{w+1}
          +\text{total factor-table size}\right),
\]
up to the cost of looking up and adding input table entries. Bit complexity
is polynomial in this quantity and the input bit length. In particular,
for every fixed $w$ this is a polynomial-time exact algorithm.
\end{theorem}
\begin{proof}
Root the bag tree and assign each factor to one containing bag. Let
$f_t(b)$ be the sum of the factors assigned to bag $t$, evaluated at its
bag assignment $b$. For a nonroot bag with parent $p(t)$, set
$D_t=B_t\cap B_{p(t)}$. Define a message indexed by separator assignments:
\[
 M_t(s)=\max_{b:\,b|_{D_t}=s}
 \left[f_t(b)+\sum_{c:\,p(c)=t}M_c(b|_{D_c})\right].
\]
Compute these messages from leaves upward, storing a maximizing bag
assignment for every separator assignment. At the root maximize the same
bracket over its full bag assignment.

We prove the recursion by induction on the subtree. The induction claim
is that $M_t(s)$ equals the maximum total of factors assigned in the
subtree rooted at $t$, among assignments of all its players agreeing with
$s$. For a leaf this is the formula. For an internal node, running
intersection implies that a player shared by two child subtrees belongs
to $B_t$, and any player shared with the parent side belongs to $D_t$.
After the bag assignment $b$ is fixed, the remaining child choices have
no conflicting unfixed player. Their maxima therefore add, and all of
them can be realized simultaneously. This proves the recursion and the
root global optimum. Follow stored choices downward to reconstruct a
consistent maximizing pure profile. Any unilateral change is among the
profiles over which global maximization was performed, so it cannot
increase the potential.

Each bag has at most $q^{w+1}$ assignments. Evaluate every bracket once,
placing its value in the corresponding separator entry and retaining the
maximum. The sum of numbers of child terms is the number of tree edges;
factor evaluations contribute at most $Lq^{w+1}$ table accesses. Every
stored value is a sum of at most $L$ rational input entries, so numerators
and denominators have polynomial bit length, even without common input
denominators. Exact comparisons and reconstruction have the asserted cost.
\end{proof}

\section{Why finding a mixed equilibrium is not a rounding argument}
\begin{proposition}\label{rounding}
The support of an exact mixed equilibrium need not consist of pure Nash
profiles. In contrast, every pure profile in the product support of a
mixed global maximizer of a multilinear potential is a global maximizer.
\end{proposition}
\begin{proof}
For the first assertion, take two players with actions $0,1$ and common
payoff $\Phi(a)=\Ind\{a_1=a_2\}$. Independent uniform mixing is an exact
equilibrium: either pure action yields expected payoff $1/2$. The supported
profile $(0,1)$ is not a pure equilibrium because either player can switch
and increase payoff from zero to one.

For the second assertion let $V$ be the largest potential at any pure
profile. A mixed potential is the average of pure potential values under
the product distribution, hence is at most $V$; a pure maximizer attains
$V$, so the mixed and pure global maxima agree. If a mixed global maximizer
put positive mass on a pure profile with value strictly below $V$, its
finite weighted average would be strictly below $V$, a contradiction.
\end{proof}
The distinction matters for attempts to turn the fixed-point construction
into a discrete algorithm. Theorem~\ref{tree} computes a global maximizer
using additional structure; the fixed-point map only ensures equilibrium.

\section{Remaining gap and checked source}
The full circuit need not expose low-width factors, and even an equivalent
expanded representation can have exponential size. Thus the structural
algorithm gives no algorithm for arbitrary succinct inputs. No hardness
reduction for $\mathsf{FIXP}\cap\mathsf{PLS}$ is supplied. The source HTML
was consulted on 9 October 2026; Question 6.1 and its footnote state the
completeness question and the solution-verification convention used here.
No computational experiments are used.
\begin{thebibliography}{9}
\bibitem{Potential} I. Anagnostides, M.-F. Balcan, K. Fragkia,
T. Sandholm, E. Tewolde, and B. H. Zhang,
\emph{The Complexity of Equilibrium Refinements in Potential Games},
arXiv:2511.03968, 2026 version. Section 3 and Question 6.1.
\url{https://arxiv.org/html/2511.03968}.
\end{thebibliography}

\section{Completion Estimate}
20\%

\end{document}
